\documentclass[12pt]{article}
\usepackage{amsmath, amssymb, amsthm}
\usepackage{geometry}
\geometry{a4paper, margin=1in}

% Theorem and Lemma Environments
\newtheorem{theorem}{Theorem}[section]
\newtheorem{lemma}[theorem]{Lemma}
\newtheorem{proposition}[theorem]{Proposition}
\newtheorem{corollary}[theorem]{Corollary}
\theoremstyle{definition}
\newtheorem{definition}[theorem]{Definition}
\newtheorem{remark}[theorem]{Remark}

\begin{document}

\title{Section 4: Non-Sobolev Compactness}
\author{}
\date{}
\maketitle

\section{Non-Sobolev Compactness}

\subsection{Compactness Framework}
To overcome limitations of classical Sobolev embeddings, we define the solution space $H$ using a scale-invariant norm:
\[
\|u\|_H = \|u\|_{L^2} + \|\nabla u\|_{L^2} + \sup_{k > k_0} \frac{E(k)}{k^\beta}, \quad \beta > 0,
\]
where:
\begin{itemize}
    \item $\|u\|_{L^2}$ is the total kinetic energy.
    \item $\|\nabla u\|_{L^2}$ measures enstrophy.
    \item $\sup_{k > k_0} \frac{E(k)}{k^\beta}$ controls high-frequency energy contributions, with $\beta$ ensuring rapid spectral decay.
\end{itemize}

This definition captures scale-invariant properties of the flow and bypasses reliance on classical Sobolev embeddings, which are often insufficient for controlling high-frequency behaviors in turbulent flows. The norm $H$ is particularly suited for flows where classical compactness arguments fail due to spectral energy cascades.

\paragraph{Comparison with Classical Sobolev Spaces:}
Unlike traditional Sobolev embeddings (e.g., $H^1$ or $H^2$), the $H$-norm explicitly incorporates high-frequency energy suppression via $\sup_{k > k_0} \frac{E(k)}{k^\beta}$. This ensures that the compactness framework remains robust even in the presence of turbulent cascades, where classical embeddings may break down.

\paragraph{Physical Interpretation:}
The $H$-norm reflects key physical properties of turbulent flows:
\begin{itemize}
    \item **Low-Frequency Components:** The terms $\|u\|_{L^2}$ and $\|\nabla u\|_{L^2}$ capture bulk energy and enstrophy, representing large-scale flow structures.
    \item **High-Frequency Suppression:** The term $\sup_{k > k_0} \frac{E(k)}{k^\beta}$ quantifies the damping of small-scale energy transfer, preventing blowup.
\end{itemize}

\subsection{Compactness Proof}
We prove strong convergence in $L^2$ for approximate solutions in $H$ without relying on classical Sobolev embeddings.

\begin{theorem}[Compactness in $H$]
Let $\{u_n\} \subset H$ be a sequence of approximate solutions satisfying:
\begin{enumerate}
    \item $\sup_n \|u_n\|_H < \infty$.
    \item $\partial_t u_n$ is uniformly bounded in $L^2(0, T; H^{-1})$.
\end{enumerate}
Then, $\{u_n\}$ has a subsequence that converges strongly in $L^2(0, T; L^2)$.
\end{theorem}

\begin{proof}
The proof adapts the Aubin--Lions lemma while leveraging spectral decay bounds instead of Sobolev embeddings. We proceed in four steps.

\paragraph{Step 1: Uniform Boundedness in $H$.}
The assumption $\sup_n \|u_n\|_H < \infty$ implies:
\begin{itemize}
    \item $\{u_n\}$ is uniformly bounded in $L^2(0, T; L^2)$.
    \item $\{\nabla u_n\}$ is uniformly bounded in $L^2(0, T; L^2)$.
    \item High-frequency contributions $\sup_{k > k_0} \frac{E_n(k)}{k^\beta}$ are uniformly controlled.
\end{itemize}
This boundedness ensures that $u_n$ belongs to the solution space $H$ for all $n$.

\paragraph{Step 2: Weak Compactness in Time.}
The boundedness of $\partial_t u_n$ in $L^2(0, T; H^{-1})$ implies that $u_n$ is weakly compact in $L^2(0, T; L^2)$. Specifically, there exists a subsequence (still denoted $u_n$) and a limit function $u \in L^2(0, T; L^2)$ such that:
\[
    u_n \rightharpoonup u \quad \text{in } L^2(0, T; L^2).
\]

\paragraph{Step 3: Spectral Decay and Approximation.}
From the Scale-Barrier Principle, the spectral energy satisfies:
\[
E(k) \leq C k^{-\gamma}, \quad \gamma > \beta + 2.
\]
This implies that for any $\epsilon > 0$, there exists a sufficiently large $k_0$ such that:
\[
\int_{k > k_0} E(k) \, dk < \epsilon.
\]
High-frequency components are thus negligible beyond a critical wave number $k_0$. Using this spectral decay, we approximate $u_n$ by smooth, compactly supported functions in Fourier space, ensuring that:
\[
\|u_n - u_n^{(k_0)}\|_{L^2} < \epsilon,
\]
where $u_n^{(k_0)}$ represents the low-pass filtered version of $u_n$ with frequencies restricted to $k \leq k_0$.

\paragraph{Step 4: Strong Convergence in $L^2$.}
Combining weak compactness from Step 2 with the approximation result from Step 3, we conclude that:
\[
\|u_n - u\|_{L^2(0, T; L^2)} \to 0 \quad \text{as } n \to \infty.
\]
The strong convergence follows from the control of high-frequency components and weak compactness in the low-frequency regime.
\qedhere
\end{proof}

\subsection{Significance}
Compactness in the space $H$ ensures smooth limits for approximate solutions and plays a critical role in establishing the existence and regularity of global solutions to the Navier--Stokes equations. Key implications include:
\begin{enumerate}
    \item **Spectral Control:** The scale-invariant norm suppresses high-frequency energy cascades, ensuring convergence without relying on classical Sobolev embeddings.
    \item **Global Regularity:** Strong convergence in $L^2$ supports the passage to smooth limits in weak formulations.
    \item **Applicability:** The framework extends naturally to fractional Sobolev norms and bounded domains, allowing broader applicability to real-world flows.
\end{enumerate}

\paragraph{Applications in Numerical Analysis:}
The compactness framework is particularly advantageous for numerical simulations. By ensuring strong convergence, the framework:
\begin{itemize}
    \item Guarantees stability in iterative schemes for solving Navier--Stokes equations.
    \item Provides error bounds for approximate solutions based on spectral decay properties.
    \item Enhances the reliability of turbulence models in capturing high-frequency dynamics.
\end{itemize}

\section*{References}
\begin{itemize}
    \item O. Ladyzhenskaya, \emph{The Mathematical Theory of Viscous Incompressible Flow}, Gordon and Breach, 1969.
    \item R. Temam, \emph{Navier--Stokes Equations: Theory and Numerical Analysis}, North-Holland, 1977.
    \item A. J. Chorin, "Spectral Methods for the Navier--Stokes Equations," J. Comp. Phys., 1968.
    \item Lions and Magenes, \emph{Non-Homogeneous Boundary Value Problems and Applications}, Springer, 1972.
    \item H. Triebel, \emph{Theory of Function Spaces}, Birkh\"auser, 1983.
    \item E. Stein, \emph{Singular Integrals and Differentiability Properties of Functions}, Princeton University Press, 1970.
\end{itemize}

\end{document}
